% Cote 119: offprints, reports, courses, CV and letters (typed and printed part): archive cote 119, site page 219 (PDF page 220).
% typescript. Draft compared: none.
% AI-transcribed by Claude Opus 5.5 (claude-opus-5-5, adaptive thinking, effort high), Message Batches API, batch r119_full (msgbatch_01DwZNm5Y7YsKBWyAN4sFs8w), 2026-11-13; not yet checked by eye.
\scanpage{119/219}
\mspage{. 6 .}

c) \uline{Théorie du groupe de Brauer} : Mes contributions
découlent pour l'essentiel de l'application de la
cohomologie étale (développée dans [16, SGA 3]) à la
théorie du groupe de Brauer. J'ai fait un exposé
d'ensemble sur les résultats connus sur ce groupe
dans [18].

\uline{Références} : M. Artin, J. Tate, J.P. Serre.

d) \uline{Théorie des algèbres de Lie} : Comme sous-produit
de recherches sur les groupes algébriques en car. $p > 0$,
je trouve certains résultats délicats sur les sous-%
algèbres de Borel ou de Cartan de certaines algèbres
de Lie, notamment sur les corps de base imparfaits
(cf. [16, SGA 6, Exp. XIII et XIV]).

\uline{Références} : M. Demazure, J. Tits, J.P. Serre.

5. \uline{Géométrie Analytique} ([10],[11],[16 bis]).

Mon influence sur la géométrie analytique
est due moins aux résultats nouveaux que j'ai pu y
démontrer (la plupart contenus dans les réf. cit.),
que par les points de vue directement inspirés par la
géométrie algébrique que j'ai pu y introduire, et les
nombreuses suggestions d'énoncés que j'ai pu y faire.
Un des plus anciens est le théorème de finitude de
Grauert pour les morphismes propres d'espaces analytiques,
aboutissant à sa généralisation récente en un théorème
qui s'énonce en termes de catégories dérivées (formula-%
tion sur laquelle j'avais insisté de longue date,
et qui a été prouvée indépendamment par R. Kiehl [$*$]
et O. Forster et K. Knorr [\unsure{$*$|$+$}]). D'autres théorèmes de

\margin{219}
